\documentclass[UTF8,fontset=fandol,11pt]{ctexart}
\usepackage[a4paper,margin=25mm]{geometry}
\usepackage{amsmath,amssymb,amsthm,mathtools}
\usepackage{xcolor}
\usepackage{booktabs,array,enumitem,fancyhdr,hyperref}
\hypersetup{colorlinks=true,linkcolor=blue!50!black,urlcolor=blue!50!black}
\setlength{\emergencystretch}{2em}
\setlength{\parskip}{0.3em}
\linespread{1.16}
\allowdisplaybreaks
\newtheorem{theorem}{定理}[section]
\newtheorem{proposition}[theorem]{命题}
\newtheorem{lemma}[theorem]{引理}
\theoremstyle{definition}
\newtheorem{definition}[theorem]{定义}
\theoremstyle{remark}
\newtheorem{remark}[theorem]{说明}
\newcommand{\dd}{\,\mathrm d}
\newcommand{\ii}{\mathrm i}
\newcommand{\Rea}{\operatorname{Re}}
\newcommand{\Ima}{\operatorname{Im}}
\newcommand{\Li}{\operatorname{Li}}
\newcommand{\supp}{\operatorname{supp}}
\pagestyle{fancy}
\fancyhf{}
\fancyhead[L]{\small 学习笔记}
\fancyhead[R]{\small 初稿}
\fancyfoot[C]{\thepage}
\setlength{\headheight}{15pt}
\author{}
\date{2026年10月22日初稿}
\makeatletter
\renewcommand{\maketitle}{%
  \thispagestyle{plain}%
  \begin{center}
    {\LARGE\bfseries \@title\par}
    \vspace{0.6em}
    {\small \@date\par}
  \end{center}
  \vspace{0.6em}}
\makeatother
\newcommand{\draftnotice}{\noindent{\small
本文为供逐篇审阅的学习笔记初稿，不是新增研究成果。
文中区分直接证明与引用结果；正文不依赖本文。
本轮已作公式、条件及排版初步核查，尚待作者逐句审定。}\par\medskip}

\title{用 Abel 求和构造 $\zeta$ 的延拓}
\begin{document}
\maketitle
\draftnotice
\section{从部分和而不是从函数方程出发}
$\zeta(s)=\sum_{n\ge1}n^{-s}$ 在 $\Rea s>1$ 绝对收敛。
若先求系数的部分和，再作分部积分，可以获得一个在更大区域内收敛的积分。
本篇构造到 $\Rea s>0$ 的亚纯延拓，并利用同一公式求 $s=1$ 附近的常数项。
这里只讨论这条构造；不把它说成建立全平面函数方程所必需的步骤。

\section{有限上限的精确恒等式}
设 $N\ge1$ 为整数。对任意复数 $s$，有限 Abel 求和给出
\[
 \sum_{n=1}^Nn^{-s}
 =N^{1-s}+s\int_1^N\lfloor x\rfloor x^{-s-1}\dd x.
\]
它也可由阶梯函数的 Stieltjes 分部积分直接得到；取 $N=1$ 时积分为零，端点项为 $1$。
写 $\lfloor x\rfloor=x-\{x\}$，对 $s\ne1$ 积分得
\[
 \sum_{n=1}^Nn^{-s}
 =\frac{s}{s-1}+\frac{N^{1-s}}{1-s}
  -s\int_1^N\{x\}x^{-s-1}\dd x.                 \tag{1}
\]
$s=1$ 时应对有限恒等式取极限或直接积分，不能单独代入两个含 $1/(s-1)$ 的项。

\section{扩大积分的收敛区域}
若 $\sigma=\Rea s>1$，式(1)的 $N^{1-s}$ 趋于零。
于是
\[
 \zeta(s)=\frac{s}{s-1}
 -s\int_1^\infty\{x\}x^{-s-1}\dd x.             \tag{2}
\]
由于 $0\le\{x\}<1$，积分的绝对值至多为
$\int_1^\infty x^{-\sigma-1}\dd x=1/\sigma$，
所以右端实际在 $\sigma>0$、$s\ne1$ 有意义。

还须说明解析性，而不能只说明逐点收敛。
任取右半平面的紧集 $K$，存在 $\delta>0$ 使 $\Rea s\ge\delta$；
参数导数的被积函数由常数倍的
$x^{-1-\delta}(1+\log x)^j$ 控制，对每个固定 $j$ 都可积。
因此可逐次求导，积分在 $\Rea s>0$ 内解析。
式(2)与原 Dirichlet 级数在公共区域一致，故构成其亚纯延拓，唯一可能的极点为 $s=1$。

\section{极点的留数与 Euler 常数}
将 $s/(s-1)=1/(s-1)+1$ 写开，积分在 $s=1$ 附近解析，立刻得到
\[
 \operatorname{Res}_{s=1}\zeta(s)=1,\qquad
 \lim_{s\to1}\left(\zeta(s)-\frac1{s-1}\right)
 =1-\int_1^\infty\frac{\{x\}}{x^2}\dd x.
\]
为识别常数，逐个单位区间计算：
\begin{align*}
 \int_n^{n+1}\frac{x-n}{x^2}\dd x
 &=\log(1+1/n)-\frac1{n+1},\\
 \int_1^N\frac{\{x\}}{x^2}\dd x
 &=\log N-\sum_{n=1}^{N-1}\frac1{n+1}
 =1+\log N-H_N.
\end{align*}
由 $\gamma_E=\lim(H_N-\log N)$，积分等于 $1-\gamma_E$，因而上述常数为 $\gamma_E$。
这里先有收敛积分，再把区间计算相加，没有对发散的两个无穷量直接相减。

\section{原级数为何仍然不能直接使用}
把式(2)中的积分拆在 $N$ 处，并与式(1)比较：
\[
 \sum_{n=1}^Nn^{-s}
 =\zeta(s)+\frac{N^{1-s}}{1-s}
   +s\int_N^\infty\{x\}x^{-s-1}\dd x,           \tag{3}
\]
对 $\Rea s>0$、$s\ne1$ 成立，且尾项模至多
$|s|N^{-\sigma}/\sigma$。

固定 $0<\sigma<1$ 时，中间项模为
$N^{1-\sigma}/|1-s|$，趋于无穷；故原部分和不收敛。
固定 $s=1+\ii t$、$t\ne0$ 时，该项为
$N^{-\ii t}/(-\ii t)$，模固定而相位不断变化。
可分别取接近 $\exp(2\pi k/|t|)$ 和 $\exp((2\pi k+\pi)/|t|)$ 的整数子列，
得到不同相位极限；取整造成的相位误差趋于零。因此部分和也不收敛。
$s=1$ 则是调和级数发散。

延拓后的函数值与原部分和不同，恰由式(3)中的显式修正项及消失尾项联系。
“函数已经延拓”并不使原来发散的级数自动变成收敛表示。

\section*{参考依据与后续核验}
所用背景为有限 Abel 求和、参数积分解析性与解析延拓唯一性，均可参看本书相关章节。
本文只构造到 $\Rea s>0$，不据此声称已经延拓到整个平面。
如需继续向左，可进一步处理周期小数部分并迭代积分；那是另一个需要完整说明的构造。
\end{document}
